% ============================================================
%  resume_standalone.tex
%  Resume independant du memoire, compilable seul.
%  Destine au rendu du 7 juin a m2imd-resps@maths.univ-amu.fr.
%  Compiler avec : pdflatex resume_standalone.tex
% ============================================================
\documentclass[11pt,a4paper]{article}

\usepackage[utf8]{inputenc}
\usepackage[T1]{fontenc}
\usepackage{lmodern}
\usepackage[french]{babel}
\usepackage[a4paper,margin=2.5cm]{geometry}
\usepackage{microtype}
\usepackage{amssymb,amsmath}
\usepackage{parskip}

\title{\vspace{-2em}Résumé de stage de Master 2 IMD}
\author{Toky Nirina RAMAHERISOA}
\date{Année universitaire 2025--2026}

\begin{document}
\maketitle
\vspace{-1.5em}

\noindent\textbf{Titre :} Génération automatique et exhaustive de
structures moléculaires bi-dimensionnelles à l'aide de la programmation
par contraintes --- focus sur les non-benzénoïdes à cycles 5, 6 et 7
atomes et validation quantique semi-empirique.

\smallskip

\noindent\textbf{Encadrants :} Nicolas PRCOVIC et Cyril TERRIOUX,
Laboratoire d'Informatique et des Systèmes (LIS), équipe COALA, Aix-Marseille
Université. \\
\noindent\textbf{Partenaires :} Yannick CARISSAN et Denis
HAGEBAUM-REIGNIER, Institut des Sciences Moléculaires de Marseille (iSm2).

\bigskip

Les hydrocarbures aromatiques polycycliques (HAP) sont des molécules
largement étudiées en chimie, pour leurs propriétés électroniques,
magnétiques et opto-électroniques. Parmi eux, les \emph{benzénoïdes}
--- HAP composés uniquement d'hexagones fusionnés par des arêtes ---
font l'objet de nombreuses études. L'étude des HAP comportant des
cycles à 5 ou 7 atomes est plus récente. La problématique centrale du
stage est la suivante : étant donné un benzénoïde, comment énumérer
toutes les substitutions admissibles, où chaque hexagone peut conserver
sa taille initiale ou être remplacé par un cycle à 5 ou 7 atomes ? La
difficulté est double. Elle est d'abord \textbf{combinatoire} : le
nombre de substitutions croît rapidement avec la taille du benzénoïde,
et seules certaines configurations sont topologiquement réalisables.
Elle est aussi liée au fait qu'\textbf{il est difficile de caractériser
avec précision la chimie} : les angles, les distances entre atomes et
la planarité --- propriété que nous cherchons précisément à évaluer
pour chaque molécule produite --- ne sont pas des grandeurs fixes, ce
qui ajoute une part de complexité au problème.

Nous abordons le problème \textbf{via} la \emph{programmation par
contraintes} (PPC), un \textbf{paradigme} déclaratif dans lequel on
décrit le problème en termes de variables et de contraintes, puis on
confie la recherche des solutions à un solveur extérieur traité comme
une \textbf{boîte noire}. Chaque molécule est représentée par un graphe
planaire, et à chaque hexagone du benzénoïde de départ est associée une
variable dont les valeurs admissibles dépendent de son voisinage local
et de la topologie globale du graphe. La PPC est ici un cadre
particulièrement adapté à la flexibilité que la chimie demande : on
peut paramétrer les contraintes, en désactiver, en ajouter, et
ré-énumérer rapidement les solutions au fil des échanges avec les
chimistes. Le paradigme a néanmoins ses limites : toutes les conditions
du problème ne se laissent pas exprimer naturellement par des
contraintes. Nous prenons donc soin d'\textbf{articuler} la résolution
par PPC avec d'autres outils. La théorie des graphes intervient pour
le pré-traitement combinatoire en amont, et la chimie computationnelle
pour la validation physique en aval, afin de tirer parti des forces de
chacun. Plusieurs variantes du modèle et différentes méthodes ont par
ailleurs été testées et comparées au cours du stage, dans une
\textbf{démarche expérimentale}.

\end{document}
